\subsection{It is enough to complete the regular part}
\label{sec:boundedregular}

The cusp and tangent cases share a substantial simplification.
We need not construct every escaping Gaussian character. It is
enough to exclude escaping contributions to the regular part,
compute that part's full rank, and construct precisely that many
local functionals. This keeps the original marker while avoiding
an unnecessary classification of the other exponential blocks.

\begin{lemma}[A positive seed on the regular part]
\label{lem:regularpositiveseed}
Let $M$ and a seed $a$ satisfy the formal Gaussian hypotheses of
Lemma~\ref{lem:seedfunctional}. After a finite cover at infinity,
write $M_0$ for the zero-exponential summand. Evaluate its dual
functionals on the canonical projection of $a$ to $M_0$. Suppose $M_0$ has a
full horizontal dual frame over $\C((\rho))$, and every member
takes $a$ into $\rho\C[[\rho]]$. Then the formal Gaussian
functional cannot take $a$ to one. No semisimplicity assumption
on the nonzero-exponential summands is required.

The same conclusion holds if the regular frame exists over
$\C((\rho))(\log\rho)$ and its seed values are finite
polynomials in $\log\rho$ with coefficients of strictly positive
$\rho$ order.
\end{lemma}
\begin{proof}
The rational seed restriction supplies a quotient $U\to\C(t)$.
The formal exponential decomposition is functorial and exact on
submodules and quotients; different exponential types have no
horizontal morphisms. These are the classical formal decomposition
facts, recalled in Sabbah~\cite[Section~3, Corollary~3.3 and its
proof]{SabbahStationary}. Consequently the quotient can only use
$U_0\subset M_0$. Since $M_0$ is trivial over the displayed field,
it is semisimple there, and this quotient extends to a horizontal
functional on $M_0$. It is a constant linear combination of its
full dual frame. Positive seed orders cannot sum to one.

Adjoining $\log\rho$, with derivative $1/\rho$, does not add
constants or negative-power exponential solutions. The same
argument applies over this differential field. A finite sum of
positive powers times logarithmic polynomials still cannot be
one. This also explains why a full regular frame, rather than
semisimplicity of the whole ambient module, is the relevant
hypothesis.
For an elementary check, expand this field in
$\C(\log\rho)((\rho))$. The operator $\rho\partial_\rho$
preserves its valuation. A nonzero exponential twist has a
negative-valuation scalar term in that operator, which cannot
be cancelled by a regular matrix. Its horizontal maps to the
trivial module therefore vanish even over this field. Its
constants are $\C$, since a rational function $a(z)$ satisfying
$a'+na=0$ is zero for integer $n\ne0$ and constant for $n=0$.
\end{proof}

\begin{proposition}[Bounded critical branches determine the regular rank]
\label{prop:boundedregularrank}
Let $P$ have degree $D\ge2$ and leading form $F_D$, and let $q$
be quadratic. Suppose, locally uniformly for small nonzero
$\epsilon$, the phase $S=P/\epsilon-q$ has a complex linear
gradient lower bound outside a compact set. The constants and
compact set may depend on $\epsilon$. Assume also
\begin{equation}\label{eq:nonisoleading}
 q_2(w)\ne0\quad\hbox{whenever }[w]\in\mathbb P^{n-1}
                           \hbox{ and }\nabla F_D(w)=0.
\end{equation}
Then the rank of the full zero-exponential summand equals the sum
of the generic local Milnor numbers of bounded phase-critical
branches whose limiting $P$ value is zero. An escaping branch
never has action polar part $\alpha/\epsilon$ for any constant
$\alpha$.
\end{proposition}
\begin{proof}
At a stationary branch the chain rule gives
$S'=-P/\epsilon^2$, hence $q=-S-\epsilon S'$.
A bounded branch has bounded $q$, so its action has only the
possible polar term $P(x(0))/\epsilon$. If a branch escapes,
divide $\nabla P=\epsilon\nabla q$ by $\|x\|^{D-1}$.
Its limiting direction is singular for $F_D$.
Condition~\eqref{eq:nonisoleading} makes the quadratic term of
$q$ along that direction nonzero. Thus $q(x(\epsilon))$ has a
pole, whereas $-S-\epsilon S'$ would be bounded for either zero
polar part or a single term $\alpha/\epsilon$. This excludes
both possibilities, without an escape-scale or Morse assumption.

We spell out the global rank comparison. Fixed-parameter tameness
gives full twisted rank equal to the total critical Milnor number
by Pham's decomposition. For the regular algebraic input
$\mathcal O$ on $\C^*_{\epsilon}\times\C^n$, form the ordinary
direct image by $(\epsilon,S)$. Its relevant cohomology is the
middle reduced Gauss--Manin module; the constant fiber cohomology
has zero exponential integral in the affine $S$ coordinate.
There is no additional finite-target discriminant from spatial
infinity. Indeed, on $S=g$ in a bounded target set,
\[
 \partial_\epsilon S=-(g+q)/\epsilon=O(1+\|x\|^2).
\]
The locally uniform gradient estimate gives a parameter-lifting
field of at most linear growth. The target-lifting field is
$O(1/\|x\|)$. Both are complete outside a compact set on each
parameter neighborhood. Thus their outer fiber pairs trivialize.
Every nonvertical discriminant curve at finite target comes
from an affine critical value, with conormal multiplicity equal
to the sum of its local Milnor numbers.

First ramify all critical-value curves into graphs. Roucairol's
nonramified rank formula, Theorem~0.1(2), Section~1.2, together
with Theorem~0.3~\cite{RoucairolFormal}, computes every nonzero
exponential-block rank from those graph multiplicities. Count
all conjugate graphs even when their polar parts agree. No
distinct-action assumption or monodromy formula is used.
Subtract these ranks from the full critical Milnor number.
Exactly the bounded-action branches remain. By the first part
they are precisely the bounded branches with limiting $P$ value
zero. This proves the assertion about the entire regular rank.
\end{proof}

We use the following explicit way to construct that regular part.
Suppose a bounded group at $P$ value $\alpha$ has coordinates,
after a cover, in which
\begin{equation}\label{eq:boundedfixedmodel}
 S=\alpha t+H_0(z)+E_\rho(z),\qquad
 \omega=\rho^\kappa J_\rho(z)\,dz,
 \quad E_\rho\in\rho\C[z][[\rho]],\quad\kappa>0,
\end{equation}
where $J_\rho$ is a unit with polynomial coefficients at each
order. If $\lambda$ is a dual functional on the fixed twisted
quotient of $H_0$, then
\[
 L(A)=e^{\alpha t}\rho^\kappa\lambda\!\left(
           e^{E_\rho}J_\rho A(x(\rho,z))\right)
\]
is a formal functional on the original module. The product rule
proves spatial exactness. Cartan's formula for the coordinate
velocity, together with the displayed moving Jacobian, proves
parameter horizontality. Every coefficient uses finitely many
polynomial terms. Fixed original markers with nonnegative powers
in these coordinates have positive seed order.

The rank check has two directions. First, every bounded critical
branch in the group must have bounded scaled coordinates $z$.
Its limiting $z$ then belongs to the critical scheme of $H_0$,
since $d_zE_\rho$ has positive order. Conversely, every isolated
critical point of $H_0$ deforms with its full local intersection
multiplicity in $H_0+E_\rho$. In the applications below the
coordinate map $x(\rho,z)$ is regular at $\rho=0$ near these
points, so all the deformed points give bounded original
critical branches. No root of the fixed model is merely counted
without being realized in the original family. Conservation of
local intersection multiplicity therefore counts exactly the
whole group, including collisions. One must also check that the
fixed critical scheme is finite and that its total Milnor number
equals the fixed twisted rank.

A polynomial basis of the fixed quotient then gives independent
functionals at leading order. Distinct groups are separated by
polynomial jets, with arbitrarily high vanishing orders to overcome
the finite negative powers introduced by the scaled representatives.
When their number equals the regular rank of
Proposition~\ref{prop:boundedregularrank}, they exhaust $M_0$.
These checks will be made explicitly below; the formal perturbation
formula alone does not establish completeness.
